% English translation of chapter7.tex lines 1806--1967.
% Source: Methods of Algebra, Volume 2, commit
% 9a5803ff77dd3257484cb177f851a73770a59dd3 (CC BY 4.0).
% stable-unit-id: o014.aljabr2.chapter7.morita-theory-b
% Coverage: Section 7.7 from the change-of-rings example to immediately
% before Lemma~\ref{prop:fg-module-cat}; continues Unit 096 and continues in Unit 098.

\hypertarget{o014.aljabr2.chapter7.morita-theory-b}{ }
% segment-id: o014.aljabr2.chapter7.morita-theory-b.e001
\begin{example}[Change of rings]\label{eg:comonad-change-ring}
	\index[sym1]{FBA@$\mathcal{F}_{B \mid A}$}
	Fix a homomorphism of $\Bbbk$-algebras $f:A\to B$. In the preceding
	framework, take $P:=B$, regarded through $f$ as an $(A,B)$-bimodule. In
	this case, the right adjoint $\Hom_{\cated{Mod}B}(B,\mathord\cdot)$ of
	$(\mathord\cdot)\dotimes{A}B$ is isomorphic, via
	$\varphi\mapsto\varphi(1)$, to the forgetful functor
	$\mathcal{F}_{B|A}:\cated{Mod}B\to\cated{Mod}A$, and
% segment-id: o014.aljabr2.chapter7.morita-theory-b.q001
	\[\begin{array}{|c|c|} \hline
		\text{unit morphism }\; \eta_N & \text{counit morphism }\; \epsilon_M \\ \hline
		N \to \mathcal{F}_{B|A}(N \dotimes{A} B) & \mathcal{F}_{B|A}(M) \dotimes{A} B \to M \\
		x \mapsto x \otimes 1 & y \otimes b \mapsto yb \\ \hline
	\end{array} \quad \begin{array}{c}
		N: \;\text{a right $A$-module,} \\
		M: \;\text{a right $B$-module.}
	\end{array}\]

	In discussions of change of rings, $\mathcal{F}_{B|A}$ is usually omitted
	from the notation for simplicity. Thus the monad $T$ and comonad $L$ may
	respectively be written as
% segment-id: o014.aljabr2.chapter7.morita-theory-b.q002
	\[\begin{array}{|c|c|c|} \hline
		T & \identity \to T & T^2 \to T \\ \hline
		N \mapsto N \dotimes{A} B & N \to N \dotimes{A} B & (N \dotimes{A} B) \dotimes{A} B \to N \dotimes{A} B \\
		& x \mapsto x \otimes 1 & (x \otimes b) \otimes b' \mapsto x \otimes bb' \\ \hline\hline
		L & L \to \identity & L \to L^2 \\ \hline
		M \mapsto M \dotimes{A} B & M \dotimes{A} B \to M & M \dotimes{A} B \to (M \dotimes{A} B) \dotimes{A} B \\
		& y \otimes b \mapsto yb & y \otimes b \mapsto (y \otimes 1) \otimes b \\ \hline
	\end{array}\]
	The required verifications are routine.
\end{example}

% segment-id: o014.aljabr2.chapter7.morita-theory-b.p001
We next consider a slightly more general situation. First, for arbitrary
right $B$-modules $P$ and $X$, regard $B$ as a $(B,B)$-bimodule in order to
define the left $B$-module
% segment-id: o014.aljabr2.chapter7.morita-theory-b.q003
\begin{equation}\label{eqn:P-vee}
	P^\vee := \Hom_{\cated{Mod}B}(P, B)
\end{equation}
and the canonical homomorphism of $\Bbbk$-modules
% segment-id: o014.aljabr2.chapter7.morita-theory-b.q004
\begin{equation}\label{eqn:dual-Hom-X}\begin{aligned}
	X \dotimes{B} P^\vee & \to \Hom_{\cated{Mod}B}(P, X) \\
	x \otimes \lambda & \mapsto x \lambda(\cdot).
\end{aligned}\end{equation}
\index[sym1]{Pvee@$P^\vee$}

% segment-id: o014.aljabr2.chapter7.morita-theory-b.p002
If $P$ is an $(A,B)$-bimodule, then $P^\vee$ has the structure of a
$(B,A)$-bimodule, $\Hom_{\cated{Mod}B}(P,X)$ has the structure of a right
$A$-module, and \eqref{eqn:dual-Hom-X} is a homomorphism of right
$A$-modules.

% segment-id: o014.aljabr2.chapter7.morita-theory-b.p003
The same construction naturally has a version with left and right
interchanged. Consequently, $P^{\vee\vee}$ is defined, and there is also a
canonical homomorphism of right $B$-modules $P\to P^{\vee\vee}$.

% segment-id: o014.aljabr2.chapter7.morita-theory-b.l001
\begin{lemma}\label{prop:AB-comonad-aux}
	Let $P$ be an $(A,B)$-bimodule which, as a right $B$-module, is finitely
	generated and projective. Then $P^\vee$, as a left $B$-module, is also
	finitely generated and projective, and in this situation:
	\begin{itemize}
		\item for every $X$, \eqref{eqn:dual-Hom-X} gives a canonical
		isomorphism of right $A$-modules;
		\item $P\to P^{\vee\vee}$ is an isomorphism.
	\end{itemize}
\end{lemma}
% segment-id: o014.aljabr2.chapter7.morita-theory-b.pf001
\begin{proof}
	As a right $B$-module, write $P$ as a direct summand of $B^{\oplus n}$ for
	some $n\in\Z_{\geq0}$. Then $P^\vee$ is naturally a direct summand of
	$B^{\oplus n}$. It is clear that replacing $P$ by $B^{\oplus n}$ in
	\eqref{eqn:dual-Hom-X} gives an isomorphism of $\Bbbk$-modules. Restricting
	to the direct summand $P$ therefore still gives an isomorphism of
	$\Bbbk$-modules, hence an isomorphism of right $A$-modules. The assertion
	$P\rightiso P^{\vee\vee}$ is proved in the same way: it is enough to check
	it for $B^{\oplus n}$.
\end{proof}

% segment-id: o014.aljabr2.chapter7.morita-theory-b.p004
Under the preceding hypothesis, the adjoint pair
\eqref{eqn:AB-adjunction} can therefore be rewritten as
% segment-id: o014.aljabr2.chapter7.morita-theory-b.q005
\begin{equation}\label{eqn:AB-adjunction-2}\begin{tikzcd}
	(\cdot) \dotimes{A} P : \cated{Mod}A \arrow[shift left, r] & \cated{Mod}B: (\cdot) \dotimes{B} P^\vee . \arrow[shift left, l]
\end{tikzcd}\end{equation}

% segment-id: o014.aljabr2.chapter7.morita-theory-b.d001
\begin{definition}\label{def:bimodule-ev-coev}
	\index[sym1]{ev-coev@$\mathrm{ev}, \mathrm{coev}$}
	Let $P$ be an $(A,B)$-bimodule which, as a right $B$-module, is finitely
	generated and projective. Define the homomorphism of $(B,B)$-bimodules
% segment-id: o014.aljabr2.chapter7.morita-theory-b.q006
	\[\begin{tikzcd}[row sep=tiny]
		\mathrm{ev}: & P^\vee \dotimes{A} P \arrow[r] & B \\
		& \lambda \otimes p \arrow[mapsto, r] & \lambda(p)
	\end{tikzcd}\]
	and the homomorphism of $(A,A)$-bimodules
% segment-id: o014.aljabr2.chapter7.morita-theory-b.q007
	\[\begin{tikzcd}[row sep=tiny]
		\mathrm{coev}: & A \arrow[r] & P \dotimes{B} P^\vee \arrow[r, "{\text{Lemma \ref{prop:AB-comonad-aux}}}" inner sep=0.6em, "\sim"'] & \End_{\cated{Mod}B}(P) \\
		& a \arrow[mapsto, rr] & & {[p \mapsto ap]} .
	\end{tikzcd}\]
\end{definition}

% segment-id: o014.aljabr2.chapter7.morita-theory-b.p005
It is easy to check that both maps are well defined. Combining
\eqref{eqn:bimod-adjunction} with Lemma~\ref{prop:AB-comonad-aux}, for every
right $A$-module $N$ and right $B$-module $M$, the unit morphism $\eta_N$ and
counit morphism $\epsilon_M$ of the adjoint pair
\eqref{eqn:AB-adjunction-2} satisfy
% segment-id: o014.aljabr2.chapter7.morita-theory-b.q008
\[\begin{tikzcd}[row sep=tiny]
	N \arrow[r, "\eta_N"] & \Hom_{\cated{Mod}B}(P, N \dotimes{A} P) & N \dotimes{A} P \dotimes{B} P^\vee \arrow[l, "\sim"'] \\
	x \arrow[mapsto, r] & {[p \mapsto x \otimes p]} & x \otimes \mathrm{coev}(1) \arrow[mapsto, l] \\
	M \dotimes{B} P^\vee \dotimes{A} P \arrow[r, "\sim"] & \Hom_{\cated{Mod}B}(P, M) \dotimes{A} P \arrow[r, "\epsilon_M"] & M \\
	x \otimes \lambda \otimes p \arrow[mapsto, r] & x \lambda(\cdot) \otimes p \arrow[mapsto, r] & x\lambda(p) \arrow[equal, d] \\
	& & (\identity_M \otimes \mathrm{ev})(x \otimes \lambda \otimes p).
\end{tikzcd}\]
Thus we may identify
% segment-id: o014.aljabr2.chapter7.morita-theory-b.q009
\begin{align*}
	\eta_N = \identity_N \otimes \mathrm{coev}: \; & N \to N \dotimes{A} P \dotimes{B} P^\vee , \\
	\epsilon_M = \identity_M \otimes \mathrm{ev}: \; & M \dotimes{B} P^\vee \dotimes{A} P \to M.
\end{align*}

% segment-id: o014.aljabr2.chapter7.morita-theory-b.p006
The preceding constructions take place at the level of the bimodules $P$ and
$P^\vee$ and do not depend on $M$ or $N$ at all.

% segment-id: o014.aljabr2.chapter7.morita-theory-b.d002
\begin{definition}\label{def:cogebroid}
	\index{module} \index{comodule}
	For an arbitrary $\Bbbk$-algebra $A$, the category
	$(A,A)\dcate{Mod}$ is monoidal under $\otimes:=\dotimes{A}$. Hence
	\S\ref{sec:alg-in-monoidal-cat} and \S\ref{sec:bialgebra} define what is
	meant by an algebra $E$ (or coalgebra $C$) in $(A,A)\dcate{Mod}$, as well
	as by left and right $E$-modules (or left and right $C$-comodules). Since
	the definition of a module (or comodule) involves $\otimes$ on only one
	side, the following slight generalization is needed.
	\begin{itemize}
% segment-id: o014.aljabr2.chapter7.morita-theory-b.i001
		\item A left $E$-module means the following data: an
		$M\in\Obj(A\dcate{Mod})$ together with a morphism in $A\dcate{Mod}$,
		$\mu_M:E\otimes M\to M$, such that the diagrams in
		Definition~\ref{def:module-algebra} commute in $A\dcate{Mod}$.

% segment-id: o014.aljabr2.chapter7.morita-theory-b.i002
		\item A left $C$-comodule means the following data: an
		$M\in\Obj(A\dcate{Mod})$ together with a morphism in $A\dcate{Mod}$,
		$\rho_M:M\to C\otimes M$, such that the diagrams in
		Definition~\ref{def:comodule-cogebra} commute in $A\dcate{Mod}$.
	\end{itemize}
	Right $E$-modules and right $C$-comodules are defined similarly, and
	morphisms between modules or comodules are defined in the standard way.
\end{definition}

% segment-id: o014.aljabr2.chapter7.morita-theory-b.p007
Applying the preceding sequence of results to describe the monad and comonad
determined by the adjoint pair \eqref{eqn:AB-adjunction-2} immediately gives
the following result.

% segment-id: o014.aljabr2.chapter7.morita-theory-b.pr001
\begin{proposition}\label{prop:Morita-comonad}
	Let $P$ be an $(A,B)$-bimodule which, as a right $B$-module, is finitely
	generated and projective, and define the $(B,A)$-bimodule $P^\vee$ as
	above. Then:
	\begin{itemize}
		\item $P\dotimes{B}P^\vee$ is an algebra in the monoidal category
		$(A,A)\dcate{Mod}$. Its multiplication is given by
% segment-id: o014.aljabr2.chapter7.morita-theory-b.q010
		\[ \identity_P \otimes \mathrm{ev} \otimes \identity_{P^\vee}: P \dotimes{B} P^\vee \dotimes{A} P \dotimes{B} P^\vee \to P \dotimes{B} P^\vee \]
		and its unit comes from
		$\mathrm{coev}:A\to P\dotimes{B}P^\vee$;
		\item $P^\vee\dotimes{A}P$ is a coalgebra in the monoidal category
		$(B,B)\dcate{Mod}$. Its comultiplication is given by
% segment-id: o014.aljabr2.chapter7.morita-theory-b.q011
		\[ \identity_{P^\vee} \otimes \mathrm{coev} \otimes \identity_P: P^\vee \dotimes{A} P \to P^\vee \dotimes{A} P \dotimes{B} P^\vee \dotimes{A} P \]
		and its counit comes from
		$\mathrm{ev}:P^\vee\dotimes{A}P\to B$.
	\end{itemize}

	Next, consider the monad $T$ on $\cated{Mod}A$ determined by the adjoint
	pair \eqref{eqn:AB-adjunction-2}, and the comonad $L$ on
	$\cated{Mod}B$ determined by it. Use Definition~\ref{def:cogebroid} when
	discussing modules and comodules.
	\begin{itemize}
		\item Specifying a $T$-module is equivalent to specifying a right
		$A$-module $N$ together with a homomorphism satisfying the usual
		properties such as associativity and unitality,
% segment-id: o014.aljabr2.chapter7.morita-theory-b.q012
		\[ N \dotimes{A} (P \dotimes{B} P^\vee) \to N; \]
		in other words, it is equivalent to making $N$ a
		$(P\dotimes{B}P^\vee)$-module.
		\item Specifying an $L$-comodule is equivalent to specifying a right
		$B$-module $M$ together with a homomorphism satisfying the usual
		properties such as coassociativity and counitality,
% segment-id: o014.aljabr2.chapter7.morita-theory-b.q013
		\[ M \to M \dotimes{B} (P^\vee \dotimes{A} P); \]
		in other words, it is equivalent to making $M$ a
		$(P^\vee\dotimes{A}P)$-comodule.
	\end{itemize}
\end{proposition}
% segment-id: o014.aljabr2.chapter7.morita-theory-b.pf002
\begin{proof}
	The proof is easier than the statement.
\end{proof}

% segment-id: o014.aljabr2.chapter7.morita-theory-b.p008
As a ring, $P\dotimes{B}P^\vee$ is simply
$\End_{B\dcate{Mod}}(P)$. By contrast, the coalgebra structure on
$P^\vee\dotimes{A}P$ looks cumbersome, but it is often more convenient in
applications. Nevertheless, comodules do not have an absolute advantage:
for example, the coalgebra must be flat for its comodule category to be
abelian. We give only an outline of a simple proof; the exercises at the end
of the chapter provide more.

% segment-id: o014.aljabr2.chapter7.morita-theory-b.pr002
\begin{proposition}\label{prop:Comod-abelian}
	Let $C$ be a coalgebra in $(B,B)\dcate{Mod}$. Write the category of right
	(or left) $C$-comodules as $\cated{Comod}C$ (or $C\dcate{Comod}$), and
	write its forgetful functor to $\cated{Mod}B$ (or $B\dcate{Mod}$) as
	$U$. If $C$, as a left (or right) $B$-module, is flat
	\cite[Definition 6.9.4]{Li1}, then $\cated{Comod}C$ (or
	$C\dcate{Comod}$) is an abelian category and $U$ is faithful and exact.
\end{proposition}
% segment-id: o014.aljabr2.chapter7.morita-theory-b.pf003
\begin{proof}
	It is enough to discuss the version for $\cated{Comod}C$. Let
	$f:M\to N$ be a homomorphism of right $C$-comodules and denote its
	cokernel at the level of $B$-modules by $\Coker(Uf)$. Then there is a
	commutative diagram in $\cated{Mod}B$ with exact rows,
% segment-id: o014.aljabr2.chapter7.morita-theory-b.q014
	\[\begin{tikzcd}
		M \arrow[r, "Uf"] \arrow[d] & N \arrow[d] \arrow[r] & \Coker(Uf) \arrow[dashed, d] \arrow[r] & 0 \\
		M \dotimes{B} C \arrow[r, "{Uf \otimes \identity}"'] & N \dotimes{B} C \arrow[r] & \Coker(Uf) \dotimes{B} C \arrow[r] & 0,
	\end{tikzcd}\]
	where the dashed arrow is uniquely determined by functoriality of the
	cokernel. It is easy to prove that this dashed arrow gives $\Coker(Uf)$ a
	comodule structure; we denote this object by $\Coker(f)$, and it is the
	cokernel of $f$. The reader may check the details if desired.

	We wish to lift $\Ker(Uf)$ to a comodule in the same way. The key is to
	ensure that
% segment-id: o014.aljabr2.chapter7.morita-theory-b.q015
	\[ 0 \to \Ker(Uf) \dotimes{B} C \to M \dotimes{B} C \xrightarrow{Uf \otimes \identity} N \dotimes{B} C \]
	is exact in $\cated{Mod}B$. This is where the flatness of $C$ as a left
	$B$-module is needed; the remaining routine verification is similar to
	the cokernel case.

	To prove that $\cated{Comod}C$ is an abelian category, it remains to show
	that the canonical morphism
% segment-id: o014.aljabr2.chapter7.morita-theory-b.q016
	\[ \Coim(f) = \Coker[\Ker(f) \hookrightarrow M] \to \Ker[N \twoheadrightarrow \Coker(f)] = \Image(f) \]
	is an isomorphism, but this need only be checked at the level of
	$\cated{Mod}B$.

	By construction, $U$ preserves kernels and cokernels and is therefore
	exact; it is plainly faithful as well.
\end{proof}

% segment-id: o014.aljabr2.chapter7.morita-theory-b.p009
Returning to the main discussion, we continue the study of the functor
$(\mathord\cdot)\dotimes{A}P$.

% Reference: Deligne, Categories tannakiennes, 4.5 Proposition.
% segment-id: o014.aljabr2.chapter7.morita-theory-b.pr003
\begin{proposition}\label{prop:Morita-comonad-ft}
	Let $P$ be an $(A,B)$-bimodule which, as a right $B$-module, is finitely
	generated and projective, and suppose that
	$(\mathord\cdot)\dotimes{A}P$ is a faithful exact functor. A right
	$A$-module $N$ is finitely generated if and only if the right $B$-module
	$N\dotimes{A}P$ is finitely generated.
\end{proposition}
% segment-id: o014.aljabr2.chapter7.morita-theory-b.pf004
\begin{proof}
	For the ``only if'' direction, suppose there are $n\in\Z_{\geq0}$ and a
	surjective homomorphism $A^{\oplus n}\twoheadrightarrow N$. Applying the
	functor gives a surjective homomorphism
	$P^{\oplus n}\twoheadrightarrow N\dotimes{A}P$, while $P$, as a right
	$B$-module, is finitely generated.

	For the ``if'' direction, take a finitely generated submodule $N'$ of $N$
	such that the image of $N'\dotimes{A}P$ contains a set of generators of
	$N\dotimes{A}P$ as a $B$-module. Since
	$(\mathord\cdot)\dotimes{A}P$ is faithful and exact, we in fact have
	$N'\dotimes{A}P\rightiso N\dotimes{A}P$, and hence $N'=N$.
\end{proof}

% segment-id: o014.aljabr2.chapter7.morita-theory-b.p010
In fact, the property that a module is finitely generated admits an entirely
categorical characterization, as follows.
